% Nouveau EGA V (cote 25): archive cote 25, site page 108 (PDF page 109).
% typescript. Draft compared: none.
% AI-transcribed by Claude Opus 5.5 (claude-opus-5-5, adaptive thinking, effort high), Message Batches API, batch ega25b (msgbatch_01A7n5YuXubUPNK1zsepi8RG), 2026-11-13; not yet checked by eye.
\scanpage{25/108}
\mspage{11}

Module \inserted{inversible} $\underline{O}_{P_i}(n)$ sur $P_i$ . Pour tout systême d'entiers
\[ \underline{n} = (n_i)_{i \in I} \in \underline{Z}^I \]
nous \xout{désignerons} poserons
\[ \underline{O}_P(\underline{n}) = \bigotimes_i \underline{O}_i(n_i) = \bigotimes_{i\,\underline{O}_S} \underline{O}_{P_i}(n_i) \quad . \]
Ceci posé, 1.1. se généralise de la façon suivante:

\uline{Proposition} 2.6. Supposons les $\underline{E}_i$ partout de rang $\geq$2. Alors pour tout
Module inversible $\underline{L}$ sur le fibré multiprojectif P, il existe une
\xout{partition} décomposition de S en somme disjointe d'ouverts $\underline{S}_{\underline{n}}$ , $\underline{n} \in \underline{Z}^I$,
et \xout{pour tout $\underline{n}$ ∈ $\underline{Z}^I$, un} un Module inversible $\underline{M}$ sur S, tels que $\underline{L}|P|S_{\underline{n}}$
soit isomorphe à $\underline{O}_P(\underline{n})|P|S_n$ . De plus, les $S_{\underline{n}}$ sont déterminés de façon
unique, et \xout{les} $\underline{M}$ est déterminé à isomorphisme unique près.

\inserted{2.7.} La démonstration \xout{résulte} consiste en une réduction immédiate à 1.1.
Sous les conditions de 2.6. , on peut donc associer à tout $\underline{L}$ inversible
sur P un ``multidegré'' $\underline{n} = (n_i)_{i \in I} \in \underline{Z}\hand{\unsure{$^{*}$}}(S)^I$,qui caractérise $\underline{L}$ à isomorphis
me près lorsque Pic(S)=0 . On peut d'ailleurs interpréter $n_i$ (appelé
``\uline{degré partiel} de $\underline{L}$ par rapport au facteur $P_i$ d'indice i '' ), lorsqu'on
se donne pour chaque i une section $g_i$ de $P_i$\xout{x} sur S (NB de telles sectio
ns existent en tous cas localement sur S), en notant que ce système définit
pour chaque i un S-morphisme \xout{$h_i$}
\[ h_i : P_i \longrightarrow P \quad ; \]
ceci posé, on a
\[ n_i = \deg h_i^{*}(\unsure{$\underline{L}$|$\underline{E}$}) \quad . \]
On fera attention qu'en général, les $n_i$ ne sont pas des entiers, mais des
applications localement constantes de S dans $\underline{Z}$ .

Procédant comme au N°1, on peut déduire de 2.6. la détermination des
morphismes d'un fibré multiprojectif dans un autre, et en particulier celle
des automorphismes d'un fibré multiprojectif. Plus intéressant pour nous,

\mspage{108}
